\scanpage{97}

% L1–2
\uline{Théorème 1} Soient $k$ un anneau, $\unsure{\mathcal E|E}$ un $k$-module, $n\in\mathbb N$ \deleted{\unsure{et des|soient}} $\unsure{\tau_i|\mathrm T_i}$ $(0\leq i\leq n)$ des pseudo réflexions dans $\mathcal E$

% L3
\[\tau_i=\mathrm{id}+a'_i\otimes a_i\qquad a'_i\in\mathcal E^\vee,\quad a_i\in\mathcal E.\tag*{(1)}\]

% L4–5
Soit $s_0\in\mathcal E$, et posons \deleted{\unsure{Cons|Puis} Posons}

% L4, right under (1)
\[(0\leq i\leq n)\]

% L6, with interline
\[\begin{aligned}s_1&=\tau_0s_0,\ s_2=\tau_1s_1\ (= \tau_1\tau_0s_0),\ \ldots,\ \inserted{s_{i+1}=\tau_i s_i\,\ldots}\\
&s_{n+1}=\tau_n s_n\ (= \tau_n\tau_{n-1}\ldots\tau_0s_0),\end{aligned}\tag*{(2)}\]

% L7, with continuation on right
\[\begin{aligned}\mathcal E_0&=ks_0,\quad\mathcal E_1=k(s_0,\deleted{\unsure{k}}a_0),\ \deleted{\unsure{\mathcal E_2=}}\ldots,\quad\mathcal E_i=k(s_0,\deleted{k}a_0,\ldots\deleted{k}a_{i-1}),\ldots.\\
&\inserted{\mathcal E_{n+1}=k(s_0,\deleted{a_0,}\ldots\deleted{k}a_n).}\end{aligned}\tag*{(3)}\]

% L8, left and below (3)
Soit enfin

% L8
\[\lambda_i=\langle s_0,a'_0\rangle\langle a_0,a'_1\rangle\ldots\langle a_{i-1},a'_i\rangle\quad(0\leq i\leq n)\tag*{(4)}\]

% L8–9, right
\margin{NB $\mathcal E_i$ eng. par $i+1$ \unsure{élts|éls}.}

% L9
\deleted{($\mathcal E_i$ engendré par $i+1$ élt). On a alors}

% L10
1°) On a

% L10–11
\[\begin{aligned}s_i&\in\mathcal E_{\unsure{j|i}}\quad\text{pour }\deleted{\unsure{i,j}}\ \inserted{0\leq i\leq j\leq n+1},\ \text{donc }\deleted{\unsure{\mathcal E}}\ k(s_0,\ldots,s_i)\subset\mathcal E_i\quad(0\leq i\leq n+1)\end{aligned}\tag*{(5)}\]
\[s_i\equiv\deleted{\unsure{s_i}}\lambda_{i-1}a_{i-1}\mod\mathcal E_{i-1}\quad(1\leq i\leq n+1)\quad\text{\uuline{si} on \uuline{suppose}}\tag*{(6)}\]

% L12, right, boxed
\[\boxed{a'_i\mid\mathcal E_{i-1}=0\quad(0\leq i\leq n)}\]

% L12–13
2°) Conditions équivalentes (A) et (B)

% L14–15
A) \inserted{On a} \deleted{Les sections $s_i$ $(0\leq i\leq n+1)$ sont lin. indépendantes,} \deleted{(3) les $(s_i)_{0\leq i\leq n+1}$ \unsure{sont|est} libre}

% L16–18
\[\tau_i\tau_j=\tau_j\tau_i\quad\text{pour }0\leq i,j\leq n,\quad i,j\text{ non consécutifs}\tag*{(\deleted{\unsure{4}}7)}\]
\[\tau_i s_0=s_0\quad\text{pour }1\leq i\leq n\tag*{(\deleted{\unsure{5}}8)}\]
\[(s_i)_{0\leq i\leq n+1}\quad\text{libre dans }\mathcal E\tag*{(\deleted{\unsure{6}}9)}\]

% L16–19, right
% slanted
\margin{il suffit que ce soit vrai \unsure{mod hom|mod hom.}, ou encore dans $\check{\mathbb P}(\mathcal E)$}

% L20
B) On a

% L21–24
\[\langle a_i,a'_j\rangle=0\quad\deleted{\unsure{chaque}}\ \text{pour }0\leq i,j\leq n,\quad i,j\text{ non conséc.}\tag*{(\deleted{\unsure{7}}10)}\]
\[\langle s_0,a'_i\rangle=0\quad\text{pour }1\leq i\leq n\tag*{(\deleted{\unsure{8}}11)}\]
\[(s_0,\deleted{\unsure{a}}a_0,\ldots,a_n)\quad\text{libre dans }\mathcal E.\tag*{(\deleted{\unsure{9}}12)}\]
\[\lambda_n\in k^*\tag*{(13)}\]

% L25
3°) De plus, ces conditions impliquent que

% L26–27
\[\mathcal E_i=k(s_0,\ldots s_{i-1})\qquad0\leq i\leq n+1\tag*{(14)}\]
\[\deleted{\unsure{(14)}\quad s_i\equiv\lambda_{i-1}a_{i-1}\quad\unsure{0\leq i}}\]

% L28–30
\[\begin{aligned}&\deleted{\mathcal E_i\ \text{stable par }\tau_0,\ldots,\tau_{i-1}\quad\unsure{(\quad<i}}\\
&\tau_j(\mathcal E_i)=\mathcal E_i\qquad j\ne i\qquad0\leq i,j\leq n\end{aligned}\tag*{(15)}\]
\[\tau_{\deleted{\unsure{j}}i}(\mathcal E_i)\ne\mathcal E_i\qquad0\leq i\leq n\qquad\text{si }k\ne0\tag*{(16)}\]

% archive sheet number: 96
